\begin{lemma}[The basis vectors of the Hermite basis: identification, coefficients, finite energy]
Fix $n \in \mathbb{N}_0$ and let $\{h_\alpha\}_{\alpha \in \mathbb{N}_0^n}$ be
the orthonormal Hermite basis of the real Hilbert space $L^2(\gamma_n)$ fixed in
\noderef{HermiteBasis}, whose vectors are the elements of \noderef{HermiteL2}.
Write $v_\alpha$ for the $\alpha$-th Hermite coefficient of
$v \in L^2(\gamma_n)$ and $\mathcal{E}$ for the fractional energy of
\noderef{Energy}.  Then the following three statements hold, each for every
multi-index $\beta \in \mathbb{N}_0^n$.
\begin{enumerate}
  \item (Identification) The $\beta$-th vector of the basis is the $\beta$-th
    Hermite element: the basis, indexed by $\mathbb{N}_0^n$, sends $\beta$ to
    $h_\beta$.
  \item (Coefficients) The Hermite coefficients of a basis vector form a
    Kronecker delta: for every $\alpha \in \mathbb{N}_0^n$,
    \[
      (h_\beta)_\alpha
        = \begin{cases} 1 & \text{if } \alpha = \beta, \\ 0 & \text{otherwise.}
          \end{cases}
    \]
  \item (Finite energy) $\mathcal{E}(h_\beta) \neq +\infty$.
\end{enumerate}
The statement contains no integral, and the only measure entering it is the one
implicit in the ambient space $L^2(\gamma_n)$; nothing further about measures is
asserted here.  The energy of clause (3) is a sum of nonnegative terms valued in
$[0,+\infty]$, so the assertion is exactly that this sum is finite.
\end{lemma}
\begin{proof}
Fix $\beta \in \mathbb{N}_0^n$ throughout; the three clauses are proved in turn,
and the proof of the third uses the second.

Clause (1).  The basis of \noderef{HermiteBasis} is by construction the
orthonormal basis assembled from the family of \noderef{HermiteL2} once that
family is known to be orthonormal, by \noderef{HermiteL2Orthonormal}, and to
have dense span, by \noderef{HermiteL2Complete}.  Unfolding that construction,
the underlying indexed family of the basis is the Hermite family itself, so
evaluating the basis at $\beta$ returns $h_\beta$.

Clause (2).  Rewriting the basis vector by clause (1), the claim becomes a
statement about the coefficients of a basis vector of an orthonormal basis.  For
such a basis the coefficient map sends the $\beta$-th basis vector to the
indicator sequence supported at $\beta$ with value $1$ there: this is the
defining property of an orthonormal basis' coordinate isomorphism applied to a
basis vector.  Reading off the $\alpha$-th entry of that sequence gives $1$ when
$\alpha = \beta$ and $0$ otherwise, which is the asserted Kronecker delta.

Clause (3).  By \noderef{Energy},
\[
  \mathcal{E}(h_\beta)
    = \sum_{\alpha \neq 0} |\alpha|^{1/2}\,\bigl|(h_\beta)_\alpha\bigr|^2 ,
\]
a sum over the multi-indices $\alpha \neq 0$ of nonnegative terms in
$[0,+\infty]$.  By clause (2) the coefficient $(h_\beta)_\alpha$ vanishes unless
$\alpha = \beta$, so
\[
  |\alpha|^{1/2}\,\bigl|(h_\beta)_\alpha\bigr|^2 = 0
  \qquad \text{for every } \alpha \neq 0 \text{ with } \alpha \neq \beta .
  \tag{Z}
\]
Now distinguish two cases, which are true for genuinely different reasons.

Case $\beta = 0$.  The summation index $\alpha$ ranges only over multi-indices
with $\alpha \neq 0$, so every $\alpha$ occurring in the sum satisfies
$\alpha \neq 0 = \beta$, and (Z) applies to all of them.  Hence every term
vanishes and the sum is a sum of zeros:
$\mathcal{E}(h_0) = 0$, which is certainly not $+\infty$.  Note that in this
case the surviving term of the other case is simply absent from the index set,
which is why the energy is literally $0$ rather than merely finite.

Case $\beta \neq 0$.  Then $\beta$ itself occurs in the index set, and (Z) says
every other term vanishes, so the sum reduces to its single term at
$\alpha = \beta$:
\[
  \mathcal{E}(h_\beta) = |\beta|^{1/2}\,\bigl|(h_\beta)_\beta\bigr|^2 .
\]
That term is the image in $[0,+\infty]$ of a real number --- a real power of the
natural number $|\beta|$ times the square of a real coefficient --- and the
canonical map from the reals into $[0,+\infty]$ never takes the value
$+\infty$.  Hence $\mathcal{E}(h_\beta) \neq +\infty$.

The two cases are exhaustive, so clause (3) holds for every $\beta$, and with it
the conjunction of the three clauses.
\end{proof}
